\documentclass[11pt,a4paper]{article}
\usepackage[T1]{fontenc}
\usepackage{lmodern}
\usepackage[margin=25mm]{geometry}
\usepackage{amsmath,amssymb,amsthm,mathtools}
\usepackage[hidelinks]{hyperref}
\usepackage{microtype}
\newcommand{\C}{\mathbb C}
\newcommand{\pd}{\operatorname{pd}}
\newcommand{\findim}{\operatorname{findim}}
\title{A special biserial algebra with finite projective dimension three}
\author{Submission for TLMC problem 00000000551}
\date{24 October 2026}
\begin{document}
\maketitle

\paragraph{Claim and scope.}
The conjecture states the universal bound $\findim\Lambda\leq2$ for special
biserial algebras. We give a finite-dimensional complex special biserial algebra
$A$ with a finitely generated module $S_3$ of projective dimension exactly $3$.
Consequently $\findim A\geq3$. The reference in the source to computations for
cyclic string algebras does not restrict its stated universal quantifier; our
example is acyclic. We do not claim a classification or an upper bound for all
special biserial algebras.

\paragraph{The algebra and its actual path presentation.}
Let $Q$ be the linearly oriented quiver
\[
3\longrightarrow2\longrightarrow1\longrightarrow0,
\]
and let $I$ kill every path of length at least two. Set $A=\C Q/I$.
Explicitly, $A=\C^4\oplus\C^3$, with
\[
(x,m)(y,n)=\left((x_i y_i)_{i=0}^3,
 (x_i n_i+m_i y_{i+1})_{i=0}^2\right).
\]
This is the standard square-zero extension with distinct left and right
diagonal actions, implemented by Mathlib's \texttt{TrivSqZeroExt}.

The formal proof also checks the path presentation, rather than presuming it.
In the incidence algebra $T$ of the four-element chain, the interval basis
$E_{ij}$, $0\leq i\leq j\leq3$, satisfies
$E_{ij}E_{kl}=\delta_{jk}E_{il}$ whenever the indices concatenate.
There is exactly one actual quiver path from $j$ to $i$ when $i\leq j$, and
none otherwise; its length is $j-i$. The formal basis and the multiplication
law therefore identify $T$ with $\C Q$. The surjective algebra homomorphism
\[
\phi:T\longrightarrow A,\qquad
f\longmapsto ((f(i,i))_{i=0}^3,(f(i,i+1))_{i=0}^2)
\]
has kernel
\[
I=\operatorname{span}_{\C}\{E_{02},E_{13},E_{03}\}
 =\langle E_{02},E_{13},E_{03}\rangle.
\]
The proof constructs the resulting quotient algebra equivalence. Thus the
relations are genuine paths of length at least two, and no nonzero vertex or
arrow relation has been introduced.

Every vertex has at most one incoming and one outgoing arrow. Every product
of two arrows lies in $I$, so each arrow has zero permitted predecessors and
successors modulo $I$. All paths of length at least two belong to $I$.
These verify the standard special biserial conditions; the algebra is also a
string algebra. The definitions used are those of Erdmann--Holm--Iyama--Schroer
\cite[pp.~168--170]{ehis}; their definition of the little finitistic dimension
is the supremum over finitely generated modules of finite projective dimension
\cite[p.~160]{ehis}.

\paragraph{Concrete projective modules and exact sequences.}
Write $S_i=\C$, with $(x,m)z=x_i z$. These are actual left $A$-modules. Each is
finitely generated, since the diagonal coordinate map $A\to S_i$ is
surjective. The module $S_0$ is a direct summand of the regular module $A$,
using the idempotent at vertex $0$, and hence is projective.

For $1\leq i\leq3$, let $P_i=\C^2$ with action
\[
(x,m)(u,v)=(x_i u,\;m_{i-1}u+x_{i-1}v).
\]
The maps $A\to P_i$, $(x,m)\mapsto(x_i,m_{i-1})$, and
$P_i\to A$, $(u,v)\mapsto(u e_i,v\alpha_{i-1})$, are $A$-linear and compose
to the identity on $P_i$. Thus each $P_i$ is a direct summand of the free
regular module, and is projective. The coordinate maps give short exact
sequences
\[
0\longrightarrow S_{i-1}\xrightarrow{v\mapsto(0,v)}
P_i\xrightarrow{(u,v)\mapsto u}S_i\longrightarrow0
\qquad(1\leq i\leq3).
\]
Injectivity, surjectivity and equality of image and kernel are checked on the
actual module maps.

The first sequence cannot split. If an $A$-linear section $s:S_1\to P_1$
existed, then $s(1)=(1,c)$ for some $c\in\C$. The arrow $\alpha_0:1\to0$
acts by zero on $S_1$, but
\[
\alpha_0s(1)=(0,1)\ne0=s(\alpha_0\cdot1),
\]
a contradiction. Hence $S_1$ is not projective.

\paragraph{Projective dimension and the contradiction.}
The three short exact sequences with projective middle terms imply
$\pd S_3\leq3$. Dimension shifting also gives
\[
\pd S_3<3\ \Longrightarrow\ \pd S_2<2
\ \Longrightarrow\ \pd S_1<1,
\]
which would make $S_1$ projective. Therefore $\pd S_3=3$.
The Lean proof uses Mathlib's \texttt{CategoryTheory.projectiveDimension},
defined through vanishing of the actual Ext groups, and its short exact
sequence dimension-shifting theorem. No dimension is assigned by definition.

The formal \texttt{finitisticDimension} is the genuine supremum of those
projective dimensions for all finitely generated objects of
\texttt{ModuleCat A} having finite projective dimension. Its indexing set
contains the explicitly constructed $S_3$, so the supremum is at least $3$.
The library's bottom value for the zero object has no effect on this positive
lower bound. This disproves the stated universal bound $2$.

\paragraph{Formal files.}
\texttt{lean/Main.lean} contains the full construction and proof, including
the path basis, quotient equivalence and special biserial certificate.
Its conclusions include
\begin{center}
\small\texttt{simple\_three\_projectiveDimension},\qquad\texttt{counterexample}.
\end{center}
The project pins Lean 4.33.0 and Mathlib commit
\begin{center}
\small\texttt{db584cd6d46c92f209a44c0f1c829460d327499d}.
\end{center}
Verification commands,
all-namespace axiom audit, fresh-kernel receipts and file hashes accompany
the submission. The complete original bilingual problem is preserved in
\texttt{SOURCE.md}.

\begin{thebibliography}{1}
\bibitem{ehis}
K.~Erdmann, T.~Holm, O.~Iyama and J.~Schroer,
\emph{Radical embeddings and representation dimension},
Advances in Mathematics \textbf{185} (2004), 159--177.
\href{https://doi.org/10.1016/S0001-8708(03)00169-5}{doi:10.1016/S0001-8708(03)00169-5}.
\end{thebibliography}
\end{document}
